\chapter{Analysis and Requirements Specification}

\section{Introduction}

This chapter formalizes the analysis phase of the RM project. It identifies system actors, specifies functional and non-functional requirements, models principal use cases, and states the project constraints that shape design and implementation decisions.

\section{Stakeholder and Actor Analysis}

\subsection{Primary Actors}

\begin{itemize}
\item \textbf{Technician}: executes interventions in the field, fills forms, attaches evidence, and updates intervention status.
\item \textbf{Supervisor}: assigns interventions, verifies execution quality, and approves or rejects completed interventions.
\item \textbf{Coordinator}: creates intervention requests, configures templates, and generates management reports.
\item \textbf{Administrator}: manages users, roles, and global system configuration.
\end{itemize}

\subsection{Secondary Stakeholders}

\begin{itemize}
\item \textbf{Client organizations}: receive validated intervention reports.
\item \textbf{Operational management}: monitors performance indicators and service quality.
\item \textbf{Technical support}: ensures platform continuity and incident resolution.
\end{itemize}

\section{Roles and Access Control}

The platform applies role-based access control (RBAC) to ensure secure and traceable operations.

\begin{table}[H]
\centering
\caption{Roles and Permissions Matrix}
\label{tab:roles_permissions_reworked}
\begin{tabularx}{\textwidth}{|l|X|X|X|X|}
\hline
\textbf{Permission} & \textbf{Technician} & \textbf{Supervisor} & \textbf{Coordinator} & \textbf{Administrator} \\
\hline
View assigned interventions & \checkmark & \checkmark & \checkmark & \checkmark \\
\hline
Update intervention status & \checkmark & \checkmark & \checkmark & \checkmark \\
\hline
Complete intervention forms & \checkmark & \checkmark & \checkmark & \checkmark \\
\hline
Assign interventions & & \checkmark & \checkmark & \checkmark \\
\hline
Approve interventions & & \checkmark & \checkmark & \checkmark \\
\hline
Create intervention requests & & & \checkmark & \checkmark \\
\hline
Configure templates & & & \checkmark & \checkmark \\
\hline
Generate reports & & & \checkmark & \checkmark \\
\hline
Manage users and roles & & & & \checkmark \\
\hline
System configuration & & & & \checkmark \\
\hline
\end{tabularx}
\end{table}

\section{Functional Requirements}

Functional requirements are grouped by capability domain.

\subsection{Authentication and Security}

\begin{enumerate}
\item The system shall authenticate users with secure credentials.
\item The system shall maintain authenticated sessions using token-based authorization.
\item The system shall enforce access control according to user roles.
\item The system shall maintain audit information on critical operations.
\end{enumerate}

\subsection{Intervention Lifecycle Management}

\begin{enumerate}
\item The system shall allow creation of intervention requests based on predefined templates.
\item The system shall support assignment of interventions to technicians.
\item The system shall support lifecycle status transitions from planning to archival.
\item The system shall allow supervisors to approve or reject intervention submissions.
\item The system shall keep a history of status changes and comments.
\end{enumerate}

\subsection{Template and Dynamic Form Management}

\begin{enumerate}
\item The system shall allow coordinators to create and manage reusable form templates.
\item The system shall generate intervention forms dynamically from selected templates.
\item The system shall validate required fields and form consistency before submission.
\end{enumerate}

\subsection{Evidence and Document Management}

\begin{enumerate}
\item The system shall allow upload and download of intervention-related files.
\item The system shall associate evidence files with interventions and users.
\item The system shall preserve file metadata for traceability.
\end{enumerate}

\subsection{Reporting and Monitoring}

\begin{enumerate}
\item The system shall generate report outputs for operational and managerial use.
\item The system shall support filtered views by project, status, and date.
\item The system shall provide dashboard information for supervision and follow-up.
\end{enumerate}

\section{Non-Functional Requirements}

\subsection{Performance}

\begin{itemize}
\item Response times should remain acceptable under concurrent usage.
\item Report generation and file operations should be optimized for operational workload.
\end{itemize}

\subsection{Reliability and Availability}

\begin{itemize}
\item The platform must ensure data persistence and recovery mechanisms.
\item Operational continuity must be maintained during normal business activity.
\end{itemize}

\subsection{Security}

\begin{itemize}
\item Authentication and authorization must protect sensitive operational data.
\item Input validation must reduce injection and data integrity risks.
\item Access to evidence and reports must be controlled by roles and permissions.
\end{itemize}

\subsection{Usability}

\begin{itemize}
\item Interfaces should remain clear for both office and field users.
\item Mobile workflows should support quick intervention data capture.
\end{itemize}

\subsection{Maintainability and Scalability}

\begin{itemize}
\item Architecture and code organization should support incremental evolution.
\item The platform should accommodate additional modules and reporting needs.
\end{itemize}

\section{System Use Cases}

\subsection{Core Use Cases by Actor}

\begin{table}[H]
\centering
\caption{Main Use Cases by Actor}
\label{tab:use_cases_by_actor}
\begin{tabularx}{\textwidth}{|l|X|}
\hline
\textbf{Actor} & \textbf{Main Use Cases} \\
\hline
Technician & Consult assigned intervention, complete dynamic form, upload evidence, submit execution output \\
\hline
Supervisor & Assign intervention, review submission, approve or reject with feedback, monitor pending items \\
\hline
Coordinator & Create intervention request, manage templates, track project progress, generate reports \\
\hline
Administrator & Manage users and roles, configure system parameters, supervise platform health \\
\hline
\end{tabularx}
\end{table}

\subsection{Representative Use Case: Intervention Validation}

\textbf{Primary actor:} Supervisor

\textbf{Preconditions:}
\begin{itemize}
\item A technician has submitted intervention data.
\item Required form fields and evidence are available.
\end{itemize}

\textbf{Main scenario:}
\begin{enumerate}
\item Supervisor opens the pending intervention.
\item Supervisor reviews captured values and attachments.
\item Supervisor either approves the intervention or rejects it with feedback.
\item System updates status history and notifies relevant actors.
\end{enumerate}

\textbf{Postconditions:}
\begin{itemize}
\item Intervention status is updated and traceable.
\item Follow-up actions become available according to the new state.
\end{itemize}

\section{Project Constraints}

\begin{itemize}
\item \textbf{Time constraint}: development and validation were performed within an academic calendar.
\item \textbf{Organizational constraint}: workflows had to align with existing operational practices.
\item \textbf{Technical constraint}: integration had to remain compatible with selected stack components and deployment model.
\item \textbf{Data constraint}: intervention data and evidence require controlled access and structured storage.
\item \textbf{Quality constraint}: solution acceptance depends on traceability, consistency, and report reliability.
\end{itemize}

\section{Chapter Conclusion}

This chapter transformed operational needs into formal requirements and constraints. It established the functional perimeter of the system and defined the quality criteria that guide architecture and implementation. The next chapter presents the design decisions adopted to satisfy these requirements.
